Volcanic activity on Io, whose rotation period synchronizes with its orbital
period, was proposed to be the result of tidal heating mainly from Jupiter.
The resultant tidal force on a body is the difference in gravitational force
experienced by the near and far sides of that body due to another body.
Measurements of the surface distortion of Io via satellite radar altimeter
mapping indicate that the surface rises and falls by up to 100\,m during
one-half orbit. Only the surface layers will move by this amount. Interior
layers within Io will move by a smaller amount, and thus we assume that on
average the entire mass of Io is moved through 50\,m. Assume that Io is
considered as two hemispheres each treated as a point mass. Calculate the
\emph{average power} of the tidal heating on Io.

\textbf{Hint:} you can use the following approximation
$(1+x)^n \approx 1 + nx$ for small $x$.

The mass of Io is $m_{\mathrm{Io}} = 8.931938 \times 10^{22}$\,kg\\
The perijove distance is $r_{peri} = 420\,000$\,km\\
The apojove distance is $r_{apo} = 423\,400$\,km\\
The orbital period of Io is 152853\,s\\
The radius of Io is $R_{\mathrm{Io}} = 1821.6$\,km.
